
\section{Tanner - History Connected Equivalency}

\subsection{Another Idea.}
Reconsider the construction of the explanation. We could add at the top an imaginary layer in which all the vertices at time $t$ are connected by arrows to all the vertices (at time $t+1$) that are adjacent to them in the Tanner graph. For the last layer, we might find that the span of a slice expands instead of shrinking. Namely, \cref{claim:prevslices} for the final slices would change to:

\begin{equation*}
  \begin{split}
\sum_{i=0}^{k}\Size{\hat{K}_{i}} +  \sum^{k}_{i=1}\Size{F_{i}} \ge \Size{\hat{K}} - \xi 
  \end{split}
\end{equation*}
Thus, following the computation at the induction stage in \cref{claim:edgesin}, the number of arrows is bounded by:

\begin{equation*}
  \begin{split}
    & \alpha \left(m - 1\right) - \beta s + d+1+ \alpha + \beta \xi \\
    & \alpha \left(m - 1\right) - \beta s + O(1) 
  \end{split}
\end{equation*}

